\section{Topology of saturated contact factorizations}
\label{sec:general-topology}
The support-orbit argument has a counterpart for arbitrary proper cones.
Here the relevant targets are the boundaries of compact cone bases.
Saturation forces these initially nonsmooth boundaries to acquire smooth
charts. This gives a joint covering theorem without a symmetry assumption
on the cones.

\subsection{Normalized cone bases and the joint covering}
Let $P,D$ be compact connected $C^1$ manifolds without boundary, both of
dimension $n\geq1$. Suppose $\zeta:P\to D$ is continuous and
\[
 s(x,z)=\sum_i\ip{A_i(x)}{B_i(z)},\qquad
 A_i:P\to K_i,\quad B_i:D\to K_i^*,
\]
where the maps are $C^1$, the cones are proper, and
$s(x,\zeta(x))=0$. Assume that the contact pairing
\begin{equation}\label{eq:abstract-contact-pairing}
 H_x(v,w)=-\sum_i\ip{DA_i(x)[v]}{DB_i(\zeta(x))[w]}
\end{equation}
is nondegenerate between $T_xP$ and $T_{\zeta(x)}D$.
No differentiability of $\zeta$ is required. Write
$c_i=\max\{\dim K_i-2,0\}$ and $I_+=\{i:c_i>0\}$.

\begin{theorem}[Joint covering at capacity equality]
\label{thm:general-joint-cover}
If $\sum_i c_i=n$, choose any $\ell_i\in\intr K_i^*$ for $i\in I_+$.
Then $A_i$ is nowhere zero, the boundary
\[
 J_i=\partial\{a\in K_i:\ell_i(a)=1\}
\]
is a $C^1$ embedded hypersurface diffeomorphic to $\Sph^{c_i}$, and
\[
 p:P\longrightarrow\prod_{i\in I_+}J_i,\qquad
 p_i(x)=\frac{A_i(x)}{\ell_i(A_i(x))},
\]
is a finite $C^1$ covering. If $P$ is simply connected, every $c_i\geq2$
and the covering has one sheet. In particular, if $P\cong\Sph^n$,
there is exactly one positive capacity, equal to $n$. If
$P\cong\prod_{a=1}^b\Sph^{p_a}$ with $p_a\geq2$, the multiset of positive
capacities is exactly $\{p_1,\ldots,p_b\}$.
\end{theorem}
\begin{proof}
At each contact every individual cone pairing is zero. Its fixed-variable
first derivatives are therefore zero:
\begin{equation}\label{eq:base-first-derivatives}
 DA_i(T_xP)\subset B_i(\zeta(x))^\perp,
 \qquad DB_i(T_{\zeta(x)}D)\subset A_i(x)^\perp.
\end{equation}
Lemma~\ref{lem:universal-channel} applies with the two tangent charts
independent. Nondegeneracy and rank subadditivity force every mixed
channel to have rank $c_i$. For $c_i>0$ this excludes both vertices,
because a two-sided derivative of a map into a pointed cone is zero at
the vertex. Complementarity puts $p_i(x)$ in $J_i$.

For $a=A_i(x)$ and $b=B_i(\zeta(x))$, the channel factors through
$b^\perp/\R a$, a space of dimension $c_i$. The quotient derivative has
rank $c_i$, since the channel does. Direct differentiation shows
\[
 Dp_i(x)[v]=0\quad\Longleftrightarrow\quad
 DA_i(x)[v]\in\R A_i(x).
\]
Thus $Dp_i$ has constant rank $c_i$. The base boundary $J_i$ is initially
only a topological sphere. The constant-rank theorem supplies a local
$c_i$-dimensional transversal on which $p_i$ is a regular embedding.
Invariance of domain makes this embedded patch relatively open in $J_i$.
Consequently $p_i(P)$ is open and compact in the connected $J_i$, and is
all of $J_i$. The regular patches give its $C^1$ embedded structure.
Radial projection from an interior point of the base is a $C^1$
diffeomorphism to $\Sph^{c_i}$: a supporting normal has strictly positive
pairing with the radial direction, so that direction is never tangent.

If $v\in\bigcap_{i\in I_+}\ker Dp_i$, the preceding kernel formula and
\eqref{eq:base-first-derivatives} give $H_x(v,w)=0$ for every $w$;
zero-capacity channels also vanish. Hence $v=0$. Source and target of
$p$ have equal dimensions, so $p$ is a local diffeomorphism. Compactness
and connectedness make it a surjective finite covering.

If $P$ is simply connected, a submersion $p_i:P\to\Sph^1$ would lift
to a real-valued function, whose maximum contradicts submersivity.
Thus the target is simply connected and the cover has one sheet.
A product of two positive-dimensional spheres has intermediate
cohomology, which proves the single-sphere assertion. More generally,
the graded indecomposable quotient $H^+/(H^+)^2$ of rational cohomology
has one generator in the dimension of each sphere. Comparing this
quotient on the source and target proves the multiset assertion.
\end{proof}

For a convex body with $C^1$ primal and polar boundaries, the independent
tangent spaces at a contact are $y^\perp$ and $x^\perp$, where
$\ip xy=1$. Their Euclidean pairing is nondegenerate: its first-slot
kernel would lie in $y^\perp\cap\R x=\{0\}$. Thus the theorem applies
without differentiating the supporting-polarity map. It also applies to
weighted product-ball contact kernels. The resulting covering need not
split as a product of maps depending on individual source coordinates.
Nor does this theorem identify an arbitrary saturated cone as Lorentz;
it identifies its normalized boundary and the capacity profile.

\begin{corollary}[Smooth ball frontier for arbitrary proper cones]
\label{cor:proper-smooth-frontier}
Let $N\geq3$ and let $d\geq3$ be an integer. Allow arbitrary proper cone factors of
dimension at most $d$. Among exact ball lifts with global $C^1$ primal
and genuine dual selections, the simultaneous minimum number of factors,
total cone dimension, and parameter of a barrier on the full ambient
cone product are
\[
 \begin{cases}
 (h,N+2h,2h),\quad h=\lceil N/(d-2)\rceil,&d<N+1,\\
 (1,N+1,2),&d\geq N+1.
 \end{cases}
\]
The ambient-parameter conclusion includes coupled barriers, and hence
also coupled logarithmically homogeneous barriers.
\end{corollary}
\begin{proof}
If $d<N+1$, every positive capacity is smaller than $N-1$.
Theorem~\ref{thm:general-joint-cover} excludes equality in the local
bound, so their sum is at least $N$. If there are $k$ such factors,
$k\geq h$ and their dimensions sum to at least $N+2k$.
Every proper cone of dimension at least two has a two-dimensional linear
section through its interior that is isomorphic to $\R_+^2$.
Restricting any ambient barrier to the product of these sections gives
an orthant barrier, whose parameter is at least $2k$
\cite[Proposition~2.3.6]{NN1994}. Extra factors cannot reduce these bounds.
For $d\geq N+1$, the local bound already gives total dimension at least
$N+1$, and the same section argument gives parameter at least two.
The grouped Lorentz lifts in Theorem~\ref{thm:smooth-lorentz}, and the
direct Lorentz lift in the second case, attain all bounds.
\end{proof}

The following local form is useful for cover arguments involving changing
active labels. It also explains why nonsmooth cone boundaries do not
remove the topological obstruction.
\begin{lemma}[Proper normalized-phase domains]\label{lem:proper-base-domain}
Let $p\geq2$, let $0<c_i<p$ with $\sum_i c_i=p$, and let $J_i$ be the
boundaries of compact convex bodies of dimensions $c_i+1$. Suppose
$U\subset\Sph^p$ is open and $q_i:U\to J_i$ are $C^1$ as ambient maps.
If the combined derivative of $q=(q_i)_i$ is injective everywhere,
then $U\ne\Sph^p$. In particular the top class of $\Sph^p$ restricts
to zero on $U$. The same holds on the inverse image of $U$ under any
coordinate projection from a product of spheres.
\end{lemma}
\begin{proof}
A map into $J_i$ has derivative rank at most $c_i$: rank $c_i+1$
would make its image contain an ambient open set. Therefore joint
injectivity forces each component to have rank $c_i$. The transversal
and invariance-of-domain argument above gives regular charts on their
images. If $U$ were the whole sphere, every image would be all of $J_i$
and $q$ would be a finite covering of their product. A circle factor
contradicts simple connectivity and finite covering degree; without
circle factors, the product is simply connected and its intermediate
cohomology contradicts that of the sphere. Finally every proper open
subset of a connected closed oriented manifold has zero ordinary top
cohomology, by noncompact-manifold duality, component by component
\cite[Section~3.3]{Hatcher2002}.
\end{proof}
